\documentclass[11pt]{article}
\usepackage[margin=1in]{geometry}
\usepackage{amsmath,amssymb}
\usepackage{booktabs}
\usepackage{url}
\title{Realistic Golfer Presentation and Lower-Limb Alignment in the Full-Body Specification (OSV-1--OSV-9, Epic \#11726, \#11731, \#11737)}
\author{UpstreamDrift engineering reference}
\date{2026-10-14}

\begin{document}
\maketitle

\noindent\textbf{Status.} Editable standalone research reference (see
\path{AGENTS.md}, Modeling Reference Documentation). It is not the
engineering design manual, whose only editable source is
\path{manuals/upstreamdrift}. Software correctness only: nothing here is a
scientific qualification.

\section{Model Topology and Assumptions}

The full-body specification (\path{full_body_spec_anthro_driver.json},
\path{full_body_spec_anthro_iron7.json}) takes its legs from the Rajagopal
OpenSense model \path{golf_humanoid.osim} through
\path{src/shared/python/motion_matching/execution/spec_builder.py}
(\texttt{leg\_extension}). Each hip is a three-coordinate joint with
primitives $R_x R_y R_z$ (flexion, adduction, rotation) and coordinates
\texttt{hip\_flexion\_\{r,l\}}, \texttt{hip\_adduction\_\{r,l\}},
\texttt{hip\_rotation\_\{r,l\}}. Knees, ankles, subtalar and
metatarsophalangeal joints are unchanged by this work.

In OpenSim, the left hip's \texttt{CustomJoint} gives
\texttt{hip\_adduction\_l} and \texttt{hip\_rotation\_l} a
\texttt{LinearFunction} coefficient of $-1$. Both hip frames have zero
orientation, so the mirroring lives only in those coefficients. The builder
read the frames but not the coefficients, so the specification's left hip
used the right hip's signs. In that convention:
\begin{itemize}
  \item $+$\texttt{hip\_rotation\_l} turned the left foot \emph{out}, while
    $+$\texttt{hip\_rotation\_r} turns the right foot in;
  \item the shared anatomical ranges (\path{range_of_motion.py}: rotation
    $\pm40^\circ$, adduction $[-50^\circ, 30^\circ]$) were reversed on the
    left side for adduction;
  \item MyoSuite, whose Rajagopal-based legs mirror the left hip, received
    the left adduction and rotation with the wrong sign (its retarget map uses
    sign $+1$).
\end{itemize}

\section{Coordinate Frames and Units}

Spec hip frames are the OpenSim hip frames post-multiplied by the permutation
$P$ (\texttt{HIP\_PERMUTATION}). This makes the base $x$ axis the OpenSim $z$
(lateral, flexion axis), base $y$ the OpenSim $x$ (anterior, adduction axis)
and base $z$ the OpenSim $y$ (superior, rotation axis). A joint maps child to
parent as
\begin{equation}
  T_{\mathrm{parent}\leftarrow\mathrm{child}}(q) =
  B \, R_x(q_1) R_y(q_2) R_z(q_3) \, F^{-1},
\end{equation}
with $B$ = \texttt{parent\_to\_base} and $F$ = \texttt{child\_to\_follower}.
Angles are in radians in documents and degrees in reports; lengths are in
metres.

\section{Equations}

Let $M = \mathrm{diag}(1,-1,-1) = R_x(\pi)$. For any angles,
\begin{equation}
  M R_x(a) M^{\top} = R_x(a), \qquad
  M R_y(b) M^{\top} = R_y(-b), \qquad
  M R_z(c) M^{\top} = R_z(-c).
\end{equation}
The mirrored hip therefore uses $B' = BM$ and $F' = FM$:
\begin{equation}
  B' R_x(a) R_y(b) R_z(c) F'^{-1}
  = B \, R_x(a) R_y(-b) R_z(-c) \, F^{-1},
\end{equation}
which is exactly OpenSim's coefficient $-1$ on adduction and rotation. At
$q = 0$ the transform is unchanged, so the zero pose, every body and every
marker are as before.

\paragraph{Zero-Twist Calibration.} \texttt{hip\_rotation\_zero} measures a
twist $\theta$ about the thigh's long axis pointing up. It is applied as
$B \leftarrow B R_z(s\theta)$, where $s = -1$ when the hip's $R_z$ axis
($F$'s third column, in the femur frame) points down the femur
($F_{2,3} < -0.5$), and $s = +1$ otherwise. A stored offset thus gives the
same femur pose with and without the mirror.

\paragraph{Symmetry Property (Tested).} With $S_p = \mathrm{diag}(1,-1,1)$ in
the pelvis \texttt{Hip} frame and $S_f = \mathrm{diag}(1,1,-1)$ in the
OpenSim femur frame, equal coordinates on both sides now give mirror-image
femurs:
\begin{equation}
  R_{\mathrm{left}}(q) = S_p \, R_{\mathrm{right}}(q) \, S_f .
\end{equation}

\section{Parameters}

No parameter is fitted or tuned. The mirrored sides are read from the
\texttt{.osim} by \texttt{spec\_builder.hip\_axis\_signs}, which returns
$\{r\!:\!+1,\ l\!:\!-1\}$ for \path{golf_humanoid.osim}. It raises an error
unless flexion has coefficient $+1$ and adduction and rotation share one sign
of $\pm1$.

\section{Numerical Method}

\begin{itemize}
  \item Builder: \texttt{leg\_extension(\ldots, mirrored\_hips)} applies
    $M$ to both hip frames of each mirrored side.
  \item Committed documents: \path{scripts/mirror_spec_left_hip.py} applies
    only $M$ to \texttt{hip\_l}, idempotently. It writes in the committed
    layout and refuses any file whose layout it would not reproduce byte for
    byte. The committed documents carry hand-set ground-support contact
    parameters that the builder does not produce, so they were not rebuilt.
  \item OpenSim full-body exports (\path{generated/full_body_anthro_*.osim})
    were regenerated from the mirrored documents with the documented exporter
    command.
\end{itemize}

\section{Acceptance Criteria (Issue \#11737)}

\begin{enumerate}
  \item The axis claim is verified against OpenSim and documented: met.
  \item The correction goes through the qualification gates, with
    before/after metrics and no tolerance change: met for the canonical
    calibrated runs (Table~\ref{tab:canon}).
  \item Both feet are within $2^\circ$ of the capture at address in five
    engines: \textbf{not met}. With the mirror alone the lead foot misses
    by $15$ to $32^\circ$ (Table~\ref{tab:feet}). With the corrections of
    Section~\ref{sec:zero} it is met in MuJoCo, OpenSim and, with the Drake
    coordinate-order fix (\#12039, PR \#12047), Drake (Table~\ref{tab:feet2}).
    Pinocchio's own FK and MyoSuite's \texttt{myolegs} FK, through the
    calibrated hip retarget of Section~\ref{sec:myo}, meet it as well
    (Table~\ref{tab:myo}).
\end{enumerate}

\section{Evidence}

\begin{table}[h]
\centering
\caption{Canonical Calibrated Runs, Marker RMS (mm): Address / IK / Dynamics}
\label{tab:canon}
\begin{tabular}{lccc}
\toprule
Run & Committed receipt & Main (control) & OSV-6 \\
\midrule
Driver (\texttt{--static-seeds}) & 7.88 / 34.05 / 84.46 & 7.22 / 33.63 / 65.00 & 7.67 / 32.98 / 61.63 \\
7-iron (\texttt{--static-seeds --zmp-filter}) & 6.56 / 31.65 / 88.61 & 6.27 / 31.65 / 52.97 & 6.55 / 31.13 / 46.97 \\
\bottomrule
\end{tabular}
\end{table}

The control re-ran the same commands on \texttt{main} at \texttt{011d6141e6}.
Most of the dynamics change since the committed receipts comes from merges
after 2026-09-27. The axis fix alone changes dynamics by $-3.4$ / $-6.0$ mm,
IK by $-0.65$ / $-0.52$ mm and address by $+0.45$ / $+0.28$ mm. There are
zero IK range-of-motion flags on both runs. The left hip-rotation range over
the swing is the exact mirror of the previous one (driver
$[-40.0, -11.4]^\circ \to [11.2, 40.0]^\circ$): the same physical motion in
the corrected sign convention.

\begin{table}[h]
\centering
\caption{Address Toe-Out (deg), Model / Capture, With \texttt{--foot-progression capture}}
\label{tab:feet}
\begin{tabular}{llcc}
\toprule
Capture & Engine & Lead (left) & Trail (right) \\
\midrule
Driver & MuJoCo & 33.1 / 16.4 & 4.1 / 4.0 \\
Driver & Drake & 31.6 / 16.4 & 4.1 / 4.0 \\
7-iron & MuJoCo & 35.9 / 15.0 & $-0.2$ / $-0.4$ \\
7-iron & Drake & 46.5 / 15.0 & $-0.2$ / $-0.4$ \\
Driver, 7-iron & OpenSim (nominal posture, FK seed) & error $< 0.005$ & error $< 0.005$ \\
\bottomrule
\end{tabular}
\end{table}

Before the fix, OSV-4 measured lead $37.0^\circ$ and trail $4.05^\circ$ on
the driver with MuJoCo. Pinocchio's full-body IK requires 41 coordinates
(the anthro document has 44), and MyoSuite is not a ground-support pipeline
plant. Neither is reported on this pathway; neither is reported as zero.

\paragraph{Owner Capture O.} Malleolus-corrected toe-out at address (33-frame
window): lead $6.1^\circ$, trail $1.6^\circ$. Its forefoot-line cross-check
reads $16.2^\circ$ / $11.7^\circ$, a disagreement of about $10^\circ$ (2 to
$5^\circ$ on the tour captures), so the headline is provisional.

\section{Address Foot Progression: Zero-Twist and Leg-Marker Corrections}
\label{sec:zero}

\paragraph{Diagnosis.} The capture stance at address is square: in the
first 24 frames of capture-A the front-waist, back-waist, knee, ankle, toe
and shoulder lines all lie within $5^\circ$ of the ankle line (capture-B
the same). Yet \texttt{hip\_rotation\_zero} returned zero twists of
$-25.7^\circ$ (right) and $+45.1^\circ$ (left) on capture-A ($-21.9$ /
$+46.7^\circ$ on capture-B). With no medial knee markers, it takes the knee
flexion axis as the normal of the plane through the hip \emph{centre} and
the \emph{lateral} knee and ankle markers. Those markers sit $w_k = 0.060$ m
and $w_a = 0.056$ m outside the joint centres along the flexion axis, which
tilts the plane. For a hip at the origin, thigh and shank of 0.40 and
0.42 m, and true zero twist, the plane normal is off by $43$, $25$ and
$17^\circ$ at 10, 20 and $30^\circ$ of knee flexion. That is the size of the
measured right-hip value, and the nearly straight lead knee explains the
larger left value. The $\pm40^\circ$ bounds of \texttt{hip\_rotation} move
with the biased zero, so the lead hip saturated.

\paragraph{Corrected Flexion Axis.} \texttt{knee\_flexion\_axis} finds the
axis $\hat n$ by the fixed-point iteration
\begin{equation}
  \hat n_{j+1} = \frac{(a - w_a \hat n_j - k + w_k \hat n_j) \times
  (k - w_k \hat n_j - h)}{\lVert\cdot\rVert},
\end{equation}
with $h$ the hip centre, $k$, $a$ the lateral markers and the cross product
ordered per side so that $\hat n$ points laterally. It stops when
$\lVert\hat n_{j+1}-\hat n_j\rVert < 10^{-12}$ (at most 50 iterations). It
recovers known twists to within $1^\circ$ at 10, 20 and $30^\circ$ of knee
flexion (unit-tested). It raises an error for a leg too straight to define
the plane.

\paragraph{Leg-Marker Calibration Constraints.} The alternating leg-marker
calibration places each marker at the prior-weighted mean of its body-frame
observations. A single lateral marker per segment lets the segment's twist
about its long axis trade against the marker's azimuth, and a sideways shift
of the forefoot pair trades against foot yaw about the ankle. On capture-A
the calibrated left forefoot line ended $28^\circ$ off square and its
midpoint 6 cm medial (about $18^\circ$ of yaw). Because the placement is a
weighted mean, the constrained least-squares placement is the orthogonal
projection of the free one. \texttt{leg\_marker\_constraints} keeps
\texttt{*KneeOut} and \texttt{*AnkleOut} on their seed azimuth (radius and
axial position free, clamped at the axis). It also gives the forefoot pair
one forward offset and pins its lateral midpoint to the seed's.

\paragraph{Scope.} Both corrections apply only with
\texttt{--foot-progression}, as the squared forefoot seeds already do. The
canonical receipts keep the legacy plane: see the failed experiment below.

\begin{table}[h]
\centering
\caption{Address Toe-Out Error (deg) and Marker RMS, \texttt{--foot-progression capture}, Lead / Trail}
\label{tab:feet2}
\begin{tabular}{llcc}
\toprule
Capture & Engine & Mirror only & With corrections \\
\midrule
Driver & MuJoCo & $+16.8$ / $+0.1$, 10.8 mm & $+0.2$ / $-0.3$, 8.3 mm \\
Driver & Drake & $+15.3$ / $+0.0$, 14.5 mm & $+0.2$ / $-0.1$, 9.5 mm \\
7-iron & MuJoCo & $+20.9$ / $+0.1$, 10.4 mm & $-0.0$ / $-1.4$, 8.1 mm \\
7-iron & Drake & $+31.5$ / $+0.2$, 11.7 mm & $-0.1$ / $+0.3$, 9.9 mm \\
\bottomrule
\end{tabular}
\end{table}

The corrected zero twists are $-5.7$ / $+25.5^\circ$ (driver) and $-1.5$ /
$+26.4^\circ$ (7-iron). Azimuth pinning alone, without midpoint pinning or
the zero-twist correction, moved the driver lead error only from $+16.8$ to
$+12.3^\circ$ and the 7-iron from $+20.9$ to $+20.2^\circ$.

\paragraph{Drake.} The Drake columns above use the coordinate-order fix of
\#12039 (PR \#12047). Before it, \texttt{DrakeFullBodyIK} passed spec-ordered
coordinates in Drake joint-creation order. The upper body (spec indices
6--29) was scrambled while the legs, which map one to one, were not: leg FK
matched MuJoCo exactly, but the centre of mass was off by 2.7 cm at the seed.
That gave lead-foot errors of $+5.7$ / $+9.7^\circ$ and address RMS of 17.4 /
9.3 mm, with \texttt{hip\_flexion\_r} pinned at $-30^\circ$.

\section{MyoSuite Leg Retarget Through the Calibrated Hip Frames (\#12052)}
\label{sec:myo}

\paragraph{Problem.} The spec-to-MyoSuite map
(\path{coordinate_map_anthro.json}) copied each leg coordinate with sign
$+1$. The pipeline fits the address in the hip-calibrated, scaled spec, where
\texttt{apply\_hip\_calibration} has rewritten each hip's
\texttt{parent\_to\_base} (functional centre, calibrated pelvis axes,
zero-twist rotation). The same \texttt{hip\_*} values then mean different
femur orientations in stock \texttt{myolegs}: the MyoSuite feet were
$-45.3$ / $+33.4^\circ$ (driver) and $-39.6$ / $+35.8^\circ$ (7-iron) off
the capture, lead / trail.

\paragraph{Frames.} Let $P = R_x(90^\circ)$ be the MyoSuite pelvis frame in
the spec pelvis (\texttt{LowerTorso} follower) frame, the \texttt{myolegs}
pelvis \texttt{body\_quat}. On the stock spec the femur body frame (follower
frame times $C^{-1}$, with $C$ the hip's \texttt{child\_to\_follower})
coincides with the MyoSuite femur frame. Hence for every hip pose,
verified to $10^{-14}$ degrees on random poses,
\begin{equation}
  R^{\mathrm{spec}}_{\mathrm{pelvis}\to\mathrm{femur}}
  = C_{\mathrm{pelvis}}^{\top}\, B\, R_x(q_1) R_y(q_2) R_z(q_3)\, C^{\top}
  = P\, R^{\mathrm{myo}}_{\mathrm{pelvis}\to\mathrm{femur}},
\end{equation}
where $B$ is the rotation of \texttt{parent\_to\_base}. MyoSuite composes
its hinges in MJCF order about femur-frame axes,
\begin{equation}
  R^{\mathrm{myo}} = R_z(f)\, R_x(s\,a)\, R_y(s\,r), \qquad
  s = +1 \text{ (right)},\ -1 \text{ (left)}.
\end{equation}
\texttt{hip\_retarget.calibrated\_hip\_targets} evaluates the spec side with
the calibrated $B$, forms $P^{\top} R^{\mathrm{spec}}$, and decomposes it as
intrinsic $Z$-$X$-$Y$ Euler angles. It refuses $|\cos a| < 10^{-6}$ (gimbal
lock). On the stock spec it returns the input coordinates to $10^{-9}$ rad
(unit-tested), so the identity map is the special case. Knee, ankle,
subtalar and mtp frames are untouched by the calibration and keep the map.

\paragraph{Left Knee Sign.} The orientation check exposed a second defect.
The spec's left knee flexes with negative \texttt{knee\_angle\_l} about the
OpenSim axis $(0, -0.071, +0.997)$, while \texttt{myolegs} flexes both knees
positive about $(0, \pm0.071, -0.997)$. The identity entry turned the left
tibia by $2|q|$ ($34.4^\circ$ at $0.3$ rad). The map entry is now $-1$. The
right spec knee uses the MyoSuite axis, so it flexes positive, yet its
range is $[-120, 10]^\circ$. Flexion is capped at $10^\circ$, and the canonical
driver IK sits on that bound. This is filed as \#12057 and is not changed
here.

\begin{table}[h]
\centering
\caption{Address Toe-Out Error (deg), Each Engine's Own FK at the Fitted
Address, \texttt{--foot-progression capture}, Lead / Trail}
\label{tab:myo}
\begin{tabular}{lccc}
\toprule
Capture & MuJoCo & Pinocchio & MyoSuite (retarget) \\
\midrule
Driver (target $16.37$ / $4.03$) & $+0.16$ / $-0.32$ & $+0.16$ / $-0.32$ & $+0.12$ / $-0.37$ \\
7-iron (target $15.00$ / $-0.37$) & $-0.00$ / $-1.41$ & $-0.00$ / $-1.41$ & $-0.04$ / $-1.46$ \\
\bottomrule
\end{tabular}
\end{table}

Before the change MyoSuite read $-45.3$ / $+33.4^\circ$ and $-39.6$ /
$+35.8^\circ$. Re-solving \texttt{hip\_rotation} against MyoSuite's FK now
moves the spec values by at most $1.5^\circ$ (driver $30.39$ / $-37.34^\circ$
against $30.25$ / $-36.95^\circ$), so the hip mapping, not the MyoSuite leg
geometry, carried the error. Femur, tibia and calcn orientations relative
to the pelvis match the calibrated spec FK within $1^\circ$ on
\texttt{myolegs}, on random leg poses for both capture specs
(\path{tests/unit/engines/myosuite/test_hip_retarget.py}). Inputs and outputs
are in \path{evidence/foot_progression/osv6_myosuite/}.

\section{Right Knee Axis Convention (\#12057)}

\subsection{Defect and Convention}
In the Rajagopal source both walker knees flex positive (range
0--120\textdegree). The primary axis is $+x$ on the right and $-x$ on the left.
\texttt{leg\_extension} built each spec knee as $R_z$ in the walker-knee frame
times \texttt{KNEE\_PERMUTATION}, and dropped the axis sign. The spec's left
knee therefore flexed negative, which matches its declared range
$[-120, 10]$\textdegree. The right knee flexed positive, so the IK could not
flex the trail knee past 10\textdegree. The spec base $z$ columns were
$(0,-0.071,-0.997)$ (right) and $(0,-0.071,+0.997)$ (left).

The chosen convention is the gait2392 one: knee flexion is negative on both
sides, i.e.\ the spec knee angle is the negated Rajagopal angle. Geometrically,
the hinge axis points to the subject's right ($+z$), so a positive $R_z$
swings the shank anteriorly. A knee whose permuted base $z$ column has a
negative $z$ component is conjugated by $A=\mathrm{diag}(1,-1,-1)$, the same
$R_x(\pi)$ used for the hip: $R_z(c)\mapsto R_z(-c)$, the identity at the
zero pose.

\begin{itemize}
\item \texttt{spec\_builder.leg\_extension} applies the convention by this
  geometric rule, independent of the OpenSim axis signs.
\item \texttt{scripts/mirror\_spec\_right\_knee.py} applies it, idempotently,
  to the committed anthropometric documents.
\item The spec-to-Rajagopal knee sign is $-1$ on both sides
  (\texttt{musculoskeletal\_graft\_validation}).
\item The MyoSuite \texttt{knee\_angle\_r} sign is $-1$, because MyoSuite
  flexes positive.
\item The MyoFullBody map already used $-1$ on both knees.
\end{itemize}

FK check (\texttt{test\_knee\_axis\_convention.py}): a $-20$\textdegree{}
right knee moves the ankle 0.148~m posterior in the pelvis frame. Before the
fix it moved the ankle 0.148~m anterior. The left knee gives 0.147~m
posterior. The zero pose is unchanged to $10^{-12}$.

\subsection{Evidence}
All runs were made on ControlTower (\texttt{ud-sim}). A control at the
parent commit \texttt{80450f267a} reproduced the previous calibrated driver
receipt exactly.

\begin{center}
\begin{tabular}{lcc}
\hline
Receipt (address / IK / dynamics, mm) & Before & After \\
\hline
Driver \texttt{--static-seeds} & 7.7 / 33.0 / 61.6 & 5.6 / 31.0 / 77.0 \\
7-iron \texttt{--static-seeds --zmp-filter} & 6.5 / 31.1 / 47.0 & 5.0 / 28.8 / 71.2 \\
Driver nominal & 43.0 / 60.5 / 119.4 & 43.7 / 56.3 / 90.4 \\
7-iron nominal (\#12030 failure) & 70.2 / 268.9 / 2743 & 70.2 / 262.9 / 3967 \\
\hline
\end{tabular}
\end{center}

\begin{itemize}
\item \textbf{Trail knee, calibrated address:}
  \begin{itemize}
  \item driver: $-0.1$ to $-26.9$\textdegree;
  \item 7-iron: $-0.6$ to $-20.8$\textdegree.
  \end{itemize}
\item \textbf{IK \texttt{knee\_angle\_r} range:}
  \begin{itemize}
  \item driver: $[-8.0, 10.0]$ to $[-51.3, -19.9]$\textdegree;
  \item 7-iron: $[-6.6, 10.0]$ to $[-49.8, -12.2]$\textdegree.
  \end{itemize}
  Both previously sat on the bound. There are no IK range-of-motion flags.
\item \textbf{Calibrated dynamics replay:} it regresses (61.6 to 77.0~mm,
  47.0 to 71.2~mm). The backswing root error grows: the driver maximum goes
  from 5.7 to 20.3~mm. The simulated finish foot slide falls from 411 to
  169~mm. The tracking controller was tuned on the capped knee.
\item \textbf{Address toe-out error, MuJoCo} (lead / trail, control at the
  parent commit to the fix,
  \path{evidence/foot_progression/knee_12057/}):
  \begin{itemize}
  \item driver: $+0.16/-0.32$ to $+0.00/-2.07$\textdegree;
  \item 7-iron: $-0.00/-1.41$ to $-0.00/-4.45$\textdegree.
  \end{itemize}
  The address marker RMS improves (8.3 to 7.5~mm and 8.1 to 7.3~mm).
\end{itemize}

With the knee flexed, femoral rotation steers the trail foot, and
\texttt{hip\_rotation\_r} sits on its $-40$\textdegree{} range limit. The trail
foot therefore misses the 2\textdegree{} OSV-6 target. The range was not
widened: that needs measured evidence (\#12109; see also \#12042). The
dynamics regression is tracked in \#12110.

\section{Hip Rewrite Pelvis Alignment (\#12109)}

\subsection{Defect}
The functional hip rewrite (\texttt{apply\_hip\_calibration}) replaces each
hip's \texttt{parent\_to\_base} $P$ by $H A_{\mathrm{new}} X$, where $H$ is
the \texttt{Hip} frame placement, $A_{\mathrm{new}}$ carries the calibrated
pelvis axes and centres, and
\begin{equation}
  X = A_{\mathrm{old}}^{-1} H^{-1} P
\end{equation}
is the hip joint frame in OpenSim pelvis coordinates. $A_{\mathrm{old}}$
must be the OpenSim-pelvis-to-\texttt{Hip} alignment that \emph{this}
document was built with. The pipeline took it from
\path{build_receipt_v2.json} for every document. That alignment is the v2
waist-marker registration, a rotation of about $50^\circ$. The
anthropometric documents are built with \texttt{pelvis\_alignment\_for}, an
axis permutation. With the foreign $A_{\mathrm{old}}$, $X$ is an arbitrary
rotation, and both calibrated hips are turned by it.

\subsection{Correction}
OpenSim's hip parent and child frames are unrotated. A builder writes the hip
as $P = H A F S$ and $C = G S$ (\texttt{child\_to\_follower}), where $F$ and
$G$ are the unrotated Rajagopal frames and $S$ is the builder's axis layout
(\texttt{HIP\_PERMUTATION}, its mirror by $M$, or the identity of the v1
layout). The rotation blocks therefore give
\begin{equation}
  A = (H^{-1} P)\, C^{\top}
\end{equation}
for each hip. \texttt{hip\_calibration.pelvis\_alignment\_from\_spec}
evaluates this for both hips and raises an error when they differ by more
than $10^{-6}$. The rewrite does not depend on the translation of $A$
(Section~\ref{sec:numerics-12109}). The function reproduces the v1 and v2
build receipts to $10^{-6}$, and returns \texttt{pelvis\_alignment\_for}
exactly for the anthropometric documents. With a calibration that reproduces
the document's own hips, the rewrite is now the identity. With the v2
receipt, it turned each anthropometric hip by more than $20^\circ$
(\path{tests/unit/motion_matching/test_pelvis_alignment_from_spec.py}).

\subsection{Numerical Note}\label{sec:numerics-12109}
With $x_t$ the translation of $X$ and $c$ the functional centre,
$A_{\mathrm{new}} = [\,R_{\mathrm{axes}} \mid c - R_{\mathrm{axes}} x_t\,]$.
The new joint origin is $H(R_{\mathrm{axes}} x_t + c - R_{\mathrm{axes}}
x_t) = Hc$ for any translation of $A_{\mathrm{old}}$. The function returns
the hip-centre midpoint as the translation.

\subsection{Evidence (Calibrated Address)}
MuJoCo, \texttt{--static-seeds --foot-progression capture}, same host
(brick), \path{evidence/foot_progression/pelvis_alignment_12109/}
(\texttt{scripts/diagnose\_address\_leg\_yaw.py}). Values are right / left.

\begin{center}
\begin{tabular}{lcccc}
\toprule
 & Driver before & Driver after & 7-iron before & 7-iron after \\
\midrule
Femur yaw at $q_{\mathrm{hip}}=0$ (deg) & 23.7 / 64.4 & $-5.7$ / 25.5 & 29.1 / 65.5 & $-1.5$ / 26.4 \\
\texttt{hip\_rotation} (deg) & $-40.0$ / 35.7 & $-17.2$ / 9.9 & $-40.0$ / 39.2 & $-23.8$ / 17.6 \\
\texttt{hip\_adduction} (deg) & $-39.6$ / 18.6 & $-6.5$ / 0.8 & $-31.6$ / 22.7 & $-4.0$ / 7.7 \\
Toe-out error, trail / lead (deg) & $-2.56$ / 0.00 & $-0.11$ / 0.22 & $-4.45$ / 0.00 & $-0.05$ / 0.19 \\
Address marker RMS (mm) & 7.38 & 7.13 & 7.29 & 4.85 \\
\bottomrule
\end{tabular}
\end{center}

After the fix, the femur yaw at zero hip coordinates equals the measured
zero-twist, and nothing else remains of the foreign rotation. Both feet are
within the $2^\circ$ OSV-6 target on both captures. No hip coordinate is on
its $\pm40^\circ$ limit. The range was not widened and no tolerance was
changed. Before the fix the solved thighs were only $-13$ to $-16^\circ$
(right) and $-1^\circ$ (left) off the pelvis forward direction. The limit
came from the rotated coordinate zero, not from the subject.

The capture's trail-foot toe-out estimators disagree by $5.1^\circ$ (driver)
and $4.0^\circ$ (7-iron) between the malleolus-corrected ankle-to-toe axis and
the forefoot normal. The $2^\circ$ criterion is tighter than that spread. The
binding target stays the corrected axis.

\section{Neutral Hip Zero-Twist and the Pelvis at Impact (\#12042)}
\label{sec:neutral-twist}

\paragraph{Problem.} The turn audit (\#12042) found the pelvis under-open at
impact: on the canonical receipts the model's hip-centre line reaches
$-40.1^\circ$ (capture-A) and $-34.5^\circ$ (capture-B) against
$-49.8^\circ$ and $-41.7^\circ$ in the capture markers. The model's own
\texttt{WaistLeft}/\texttt{WaistRight} sites, evaluated with the receipt's
calibrated attachments, reach only $-42.6^\circ$ on capture-A. So the gap is a
fitting deficit, not a difference in the line definition. At impact
\texttt{hip\_rotation\_l} sits on its $+40^\circ$ bound on both captures.

\paragraph{Mechanism.} \texttt{hip\_rotation\_zero} estimates the thigh twist
$\tau$ over the first 24 frames, and \texttt{apply\_hip\_calibration} rotates
each femur frame about its long axis by $\tau$. The coordinate zero is
therefore the measured address twist, and the fixed $\pm40^\circ$ range of
\texttt{hip\_rotation} (the lower-limb ranges of
\texttt{range\_of\_motion}) becomes a window shifted by $\tau$ in
pelvis-relative terms. Without the
lateral-marker correction (Section~\ref{sec:zero}), the canonical
\texttt{--static-seeds} specs carry $\tau = +45.1 / -25.7^\circ$ (lead /
trail, capture-A) and $+46.7 / -21.9^\circ$ (capture-B). The address stance
is square, so these are estimation errors, not posture. Even the corrected
plane leaves $+25.5^\circ$ on the nearly straight lead knee. The lead window
is shifted by tens of degrees, and the internal rotation the impact pose needs
falls outside it.

\paragraph{Option.} \texttt{--hip-zero-twist neutral}
(\texttt{hip\_calibration.select\_hip\_zero\_twist}) applies $\tau = 0 / 0$:
the femur frame is aligned with the calibrated pelvis frame (the hip-rewrite
alignment of Section~\ref{sec:numerics-12109}). The $\pm40^\circ$ bound is
then anatomical, and the address twist is solved as a coordinate value. The
default \texttt{capture} reproduces the canonical receipts. Each receipt
records \texttt{hip\_zero\_twist\_mode} and the measured twist
(\texttt{hip\_zero\_twist\_measured\_deg}) next to the applied one. No bound
or tolerance changes.

\paragraph{Evidence.} MuJoCo, ControlTower \texttt{ud-sim}, canonical
commands (driver \texttt{--static-seeds}, 7-iron \texttt{--static-seeds
--zmp-filter}), same commit. The same-commit baselines reproduce the canonical
receipts exactly. Turn errors are model minus markers (top / impact, degrees).
FD is the whole-swing marker RMS, with address-to-impact / after-impact in
brackets. Receipts: \path{docs/development/full_body_models/evidence/hip_zero_twist_12042/}.

\begin{center}
\footnotesize
\setlength{\tabcolsep}{3pt}
\begin{tabular}{lcccccc}
\toprule
Run & Pelvis & Upper & Shoulder & Lead hip rot. & Address / IK & FD \\
 & & trunk & girdle & at impact & (mm) & (mm) \\
\midrule
Driver, capture & $-0.5$ / $+9.7$ & $+18.3$ / $+4.0$ & $+14.6$ / $-5.5$ & $40.0$ & 6.3 / 29.5 & 80.8 (34.2 / 145.0) \\
Driver, neutral & $+3.2$ / $+2.5$ & $+17.7$ / $+5.6$ & $+12.0$ / $-9.0$ & $27.0$ & 5.3 / 25.6 & 51.2 (31.6 / 83.8) \\
7-iron, capture & $+2.7$ / $+7.2$ & $+13.6$ / $+6.1$ & $+10.8$ / $+3.3$ & $40.0$ & 6.6 / 26.8 & 54.5 (30.1 / 92.1) \\
7-iron, neutral & $+1.4$ / $+2.8$ & $+13.9$ / $+5.2$ & $+11.0$ / $+8.3$ & $27.7$ & 2.1 / 22.9 & 91.2 (34.0 / 165.3) \\
\bottomrule
\end{tabular}
\end{center}

With the neutral twist, the pelvis line is within $3.2^\circ$ of the markers
at top and impact on both captures (the \#12042 bound is $5^\circ$). The lead
hip leaves its bound at impact. The lead-leg IK residual falls from 36 to
15~mm (capture-A) and from 40 to 16~mm (capture-B). The address marker RMS
improves on both captures. Neither run has a range-of-motion flag.

\paragraph{Limitations.}
\begin{itemize}
  \item The forward-dynamics replay is mixed. Capture-A improves at every
    phase. On capture-B the address-to-impact FD rises by 3.9~mm, near the
    about 3~mm noise floor measured for \#12042 slice 7, while after-impact FD
    rises from 92 to 165~mm. That is the follow-through rollout sensitivity
    already documented for \#12042.
  \item The trail \texttt{hip\_rotation\_r} now reaches its $-40^\circ$ bound
    late in the downswing on both captures ($-39.5$ / $-38.3^\circ$ at
    impact).
  \item The upper-trunk over-twist at the top ($+14$ to $+18^\circ$) is
    unchanged, as is the shoulder-girdle line at impact. The neutral twist
    does not address those (slice 7).
  \item The neutral setting assumes zero femoral twist relative to the pelvis
    frame at the coordinate zero. It does not model individual femoral version.
    A functional knee axis over the swing remains the alternative estimate.
  \item This is inverse kinematics plus feedback-assisted tracking, not an
    open-loop dynamics result.
\end{itemize}

\paragraph{Failed Experiment: Widening the Hip Rotation Range.} With the
range opened to $\pm60^\circ$ (an experiment-only override, never proposed),
the capture-A IK runs straight to the new bound ($59.4^\circ$ at impact). On
capture-B it lands on a different leg branch: the pelvis at impact gets worse
($+10.5^\circ$) and FD rises to 120.7~mm. The range is not the cause; its
placement is.

\paragraph{Reproduce.}
\begin{verbatim}
python3 -m src.shared.python.motion_matching.pipeline.cli \
  --spec docs/development/full_body_models/full_body_spec_anthro_driver.json \
  --capture driver --static-seeds --hip-zero-twist neutral --out runs/hzn_driver
python3 -m src.shared.python.motion_matching.pipeline.cli \
  --spec docs/development/full_body_models/full_body_spec_anthro_iron7.json \
  --capture iron --static-seeds --zmp-filter --hip-zero-twist neutral \
  --out runs/hzn_iron
\end{verbatim}

\section{Two-Hand Golf Grip Definition and Kinetics (OSV-2, OSV-7, OSV-9)}
\label{sec:grip}

\paragraph{Coordinate Frames and Geometry.}
Grip geometry is defined in the club coordinate frame $\mathcal{F}_c = (\mathbf{p}_c, \hat{\mathbf{x}}_c, \hat{\mathbf{y}}_c, \hat{\mathbf{z}}_c)$, where $\hat{\mathbf{z}}_c$ is the longitudinal shaft axis directed from the clubhead toward the grip butt cap, $\hat{\mathbf{x}}_c$ is the toe axis orthogonal to the shaft in the face-plane direction, and $\hat{\mathbf{y}}_c = \hat{\mathbf{z}}_c \times \hat{\mathbf{x}}_c$ is the face-normal axis pointing anteriorly from the clubface.

The lead hand (left hand for a right-handed golfer) is placed adjacent to the butt cap with an offset $\Delta z_{\mathrm{lead}} = 15\text{ mm}$:
\begin{equation}
  \mathbf{p}_{\mathrm{lead}} = \mathbf{p}_{\mathrm{butt}} - \Delta z_{\mathrm{lead}} \hat{\mathbf{z}}_c .
\end{equation}
The trail hand (right hand) is placed immediately distal to the lead hand along the shaft axis:
\begin{equation}
  \mathbf{p}_{\mathrm{trail}} = \mathbf{p}_{\mathrm{lead}} - L_{\mathrm{sep}} \hat{\mathbf{z}}_c ,
\end{equation}
where $L_{\mathrm{sep}} = \|\mathbf{p}_{\mathrm{trail}} - \mathbf{p}_{\mathrm{lead}}\|$ is the hand centre separation distance reconciled from the full-body specification closure definitions (\texttt{closure.placement\_a} and \texttt{placement\_b}).

\paragraph{Neutral Grip Orientations and V-Angles.}
Anatomical grip closure specifies:
\begin{enumerate}
  \item \textbf{Lead Hand}: The ``V'' formed by the thumb and index finger web points in the quadrant between the golfer's chin and trail shoulder ($\theta_{V,\mathrm{lead}} \in [15^\circ, 25^\circ]$). The thumb sits slightly trail-side of the top-dead-centre line of the grip (short thumb).
  \item \textbf{Trail Hand}: The thenar eminence (life line) wraps conformally over the dorsal aspect of the lead thumb. The trail ``V'' aligns parallel to the lead ``V'' ($\mathbf{v}_{\mathrm{trail}} \parallel \mathbf{v}_{\mathrm{lead}}$).
  \item \textbf{Closure Style}: The default configuration is the Vardon overlap grip, with optional interlock and ten-finger variants supported by the geometry generator (\texttt{src/shared/python/model\_appearance/grip\_pose.py}).
\end{enumerate}

\paragraph{Kinematic Closure and Compliant Kinetics.}
In OpenSim (\texttt{golf\_humanoid.osim} and musculoskeletal exports), both hands are constrained to the club. Kinematic loop closure enforces:
\begin{equation}
  \mathbf{\Phi}_{\mathrm{grip}}(\mathbf{q}) = 
  \begin{bmatrix}
    \mathbf{p}_{\mathrm{hand\_l}}(\mathbf{q}) - \mathbf{p}_{\mathrm{grip\_lead}}(\mathbf{q}) \\
    \mathbf{p}_{\mathrm{hand\_r}}(\mathbf{q}) - \mathbf{p}_{\mathrm{grip\_trail}}(\mathbf{q})
  \end{bmatrix} = \mathbf{0},
\end{equation}
with the closure residual bounded by $\|\mathbf{\Phi}_{\mathrm{grip}}(\mathbf{q}(t))\| \le 5\text{ mm}$ throughout the entire swing. For compliant kinetic parity (OSV-7), distributed bushing elements apply restorative wrenches:
\begin{equation}
  \mathbf{W}_{\mathrm{bushing}} = \mathbf{K}_b \Delta \mathbf{x} + \mathbf{D}_b \Delta \mathbf{v},
\end{equation}
matching the multi-engine KKT multiplier contact models across MuJoCo, Drake, Pinocchio, and OpenSim.

\section{Head Gaze Stabilisation Schedule and Neck IK (OSV-3)}
\label{sec:gaze}

\paragraph{Coordinate Frames and Gaze Metric.}
In the world frame $\mathcal{W}$ ($Z$-up, golfer facing $-X$, target line $-Y$):
\begin{itemize}
  \item \textbf{Eye Midpoint}: $\mathbf{e}(\mathbf{q}) = \mathbf{p}_{\mathrm{head}}(\mathbf{q}) + \mathbf{R}_{\mathrm{head}}(\mathbf{q}) \mathbf{d}_{\mathrm{eye}}$, using the de Leva (1996) anthropometric head midpoint offset $\mathbf{d}_{\mathrm{eye}} = [0.065, 0.0, 0.040]^\top\text{ m}$.
  \item \textbf{Ball Location at Address}: $\mathbf{b} \in \mathbb{R}^3$, established at the clubface centre at address offset by one ball radius $r_{\mathrm{ball}} = 21.335\text{ mm}$ along the face normal, resting on the ground plane ($z = r_{\mathrm{ball}}$).
  \item \textbf{Gaze Direction}: $\mathbf{g}(t) = \mathbf{R}_{\mathrm{head}}(t) \hat{\mathbf{g}}_{\mathrm{head}}$, where $\hat{\mathbf{g}}_{\mathrm{head}} = [1, 0, 0]^\top$ is the anatomical forward facial axis.
  \item \textbf{Angular Gaze Error}:
    \begin{equation}
      \theta_{\mathrm{gaze}}(t) = \arccos\left( \frac{\mathbf{g}(t) \cdot (\mathbf{b} - \mathbf{e}(t))}{\|\mathbf{b} - \mathbf{e}(t)\|} \right) .
    \end{equation}
\end{itemize}

\paragraph{Time-Dependent Gaze Schedule.}
To provide human-like head stabilisation, the gaze target $\mathbf{r}^*(t)$ is scheduled over three distinct swing phases:
\begin{equation}
  \mathbf{r}^*(t) = 
  \begin{cases}
    \mathbf{b}, & t \le t_{\mathrm{impact}} + t_{\mathrm{hold}} \\
    (1 - s(\tau)) \mathbf{b} + s(\tau) \left( \mathbf{e}(t) + d_{\mathrm{target}} \hat{\mathbf{x}}_t \right), & t_{\mathrm{impact}} + t_{\mathrm{hold}} < t < t_{\mathrm{impact}} + t_{\mathrm{hold}} + T_{\mathrm{rel}} \\
    \mathbf{e}(t) + d_{\mathrm{target}} \hat{\mathbf{x}}_t, & t \ge t_{\mathrm{impact}} + t_{\mathrm{hold}} + T_{\mathrm{rel}}
  \end{cases}
\end{equation}
where $t_{\mathrm{hold}} = 0.03\text{ s}$ (30 ms gaze hold through impact), $T_{\mathrm{rel}} = 0.35\text{ s}$ (release duration), $\hat{\mathbf{x}}_t = -\hat{\mathbf{y}}$ (down target line), and $s(\tau) = 10\tau^3 - 15\tau^4 + 6\tau^5$ is the quintic minimum-jerk transition function of normalised release time:
\begin{equation}
  \tau = \frac{t - (t_{\mathrm{impact}} + t_{\mathrm{hold}})}{T_{\mathrm{rel}}} \in [0, 1] .
\end{equation}

\paragraph{Neck Inverse Kinematics.}
Head stabilization is solved through the neck degrees of freedom (\texttt{NeckInputX/Y/Z}, ROM: flexion $\pm 45^\circ$, lateral tilt $\pm 60^\circ$, axial rotation $\pm 80^\circ$) by adding a soft regularisation penalty to the tracking objective:
\begin{equation}
  J_{\mathrm{fit}}(\mathbf{q}) = J_{\mathrm{marker}}(\mathbf{q}) + w_{\mathrm{gaze}} \theta_{\mathrm{gaze}}^2(\mathbf{q}) .
\end{equation}
This regularisation is configurable and strictly receipt-tracked; marker-faithful fits remain available without alteration.

\section{Neutral Address Foot Orientation and Progression Parity (OSV-4, OSV-6)}
\label{sec:foot_progression_osv}

\paragraph{Coordinate System and Foot Long Axis.}
The address reference frame is oriented with up axis $\hat{\mathbf{z}}$, target axis $\hat{\mathbf{x}}_t = -\hat{\mathbf{y}}$, and golfer-forward axis $\hat{\mathbf{f}} = \hat{\mathbf{z}} \times \hat{\mathbf{x}}_t = -\hat{\mathbf{x}}$ directed toward the ball line.

Using ground-projected markers:
\begin{align}
  \mathbf{p}_{\mathrm{heel}} &= \mathbf{p}_{\mathrm{AnkleOut}} - (\mathbf{p}_{\mathrm{AnkleOut}} \cdot \hat{\mathbf{z}}) \hat{\mathbf{z}} , \\
  \mathbf{p}_{\mathrm{toe}} &= \frac{1}{2}(\mathbf{p}_{\mathrm{ToeIn}} + \mathbf{p}_{\mathrm{ToeOut}}) - \left( \frac{1}{2}(\mathbf{p}_{\mathrm{ToeIn}} + \mathbf{p}_{\mathrm{ToeOut}}) \cdot \hat{\mathbf{z}} \right) \hat{\mathbf{z}} .
\end{align}
The longitudinal foot unit vector is $\hat{\mathbf{a}} = \frac{\mathbf{p}_{\mathrm{toe}} - \mathbf{p}_{\mathrm{heel}}}{\|\mathbf{p}_{\mathrm{toe}} - \mathbf{p}_{\mathrm{heel}}\|}$.

\paragraph{Foot Progression Angle.}
The signed foot progression (toe-out) angle is computed as:
\begin{equation}
  \theta_{\mathrm{prog}} = \mathrm{atan2}\left( \hat{\mathbf{z}} \cdot (\hat{\mathbf{f}} \times \hat{\mathbf{a}}),\, \hat{\mathbf{f}} \cdot \hat{\mathbf{a}} \right) ,
\end{equation}
where positive $\theta_{\mathrm{prog}}$ designates outward flare (lead foot rotated toward target line, trail foot rotated away from target line). The address stance value is the median over the initial static stance window prior to takeaway. When markers are absent or uncalibrated, a default neutral stance of $\theta_{\mathrm{prog}} = 20^\circ$ per foot is applied.

\paragraph{Interaction With Hip Mirroring (OSV-6).}
Foot progression at address is achieved through the hip rotation coordinates (\texttt{hip\_rotation\_r}, \texttt{hip\_rotation\_l}) in conjunction with subtalar orientation. OpenSim's negative coefficient ($-1$) on \texttt{hip\_rotation\_l} ensures that positive lead hip rotation produces anatomical external rotation (toe-out). Prior to the hip mirroring fix, unmirrored kinematics forced the lead hip into internal rotation limit ($40^\circ$), causing severe lead-foot toe-in and preventing natural address posture.

\section{Limitations and Failed Experiments}

\begin{itemize}
  \item \textbf{MyoSuite retarget scope.} The evidence is kinematic: the
    \texttt{myolegs} pose at the fitted address, with the knee's
    equality-coupled helper joints at 0 (in simulation they follow the knee
    angle and add up to several degrees of tibial rotation). The pelvis is
    placed by the evaluation script, because the map has no pelvis
    orientation entry. Only orientations are matched: the calibrated hip
    centre is not transferred, and \texttt{replay.ReplayConfig} needs the fit's
    spec (\texttt{hip\_spec}) to use the retarget. It is not a MyoSuite
    dynamics or muscle result.
  \item \textbf{Default-on zero-twist correction (failed).} Applied to the
    canonical \texttt{--static-seeds} runs, the correction gave driver
    address / IK / dynamics of 7.35 / 36.81 / 55.06 mm (main control 7.22 /
    33.63 / 65.00). The 7-iron gave 6.59 / 31.82 / 65.34 mm (main 6.27 /
    31.65 / 52.97). It also cut the driver reference finish ZMP-inside
    fraction from 0.272 to 0.061, with \texttt{hip\_rotation} reaching its
    $\pm40^\circ$ bounds through the swing. The remaining left zero twist
    of about $25^\circ$ still reflects the ill-conditioned plane of a nearly
    straight lead knee. A functional knee axis estimated over the swing is
    the open alternative.
  \item At address the solved hip rotations still shared a common component
    (7-iron $-40$ / $+36^\circ$, driver $-37$ / $+30^\circ$, right / left).
    \#12109 traced it to the foreign pelvis alignment in the hip rewrite
    (Section above), not to the zero-twist or the pelvis yaw. The left
    zero-twist of about $25^\circ$ remains and is still the open item.
  \item With the mirror alone (Table~\ref{tab:feet}), the lead foot is too
    far turned out: \texttt{hip\_rotation\_l} is pinned at its $40^\circ$
    internal limit in every engine run.
  \item The legacy documents \path{full_body_spec_v1.json} and
    \path{full_body_spec_v2.json}, built by
    \path{docs/development/full_body_models/build_full_body_spec.py} (a
    separate copy of the builder), keep the unmirrored left hip.
  \item Historical receipts that cite the old document hashes (grip
    kinetics, \texttt{ms102} structural receipts, the \texttt{downswing\_*}
    experiments under \texttt{anthro\_driver}) are history and were not
    regenerated.
  \item This is inverse kinematics plus feedback-assisted tracking. It is
    not an independently replayed open-loop dynamics result.
  \item \textbf{Grip mesh and contact fidelity (OSV-2, OSV-7).} Curled hand meshes
    use rigid visual geometry baked at the specified grip radius; live soft-tissue
    pad deformation and individual finger contact mechanics are not solved in IK.
    Distributed bushing kinetics provide multi-engine kinetic stability.
  \item \textbf{Gaze regularisation scope (OSV-3).} Head gaze tracking applies
    a soft regularisation term in neck inverse kinematics; head translation is not
    artificially locked in world space and follows torso/cervical kinematics.
  \item \textbf{Foot progression default (OSV-4).} The default $20^\circ$ toe-out
    per foot represents a captured anatomical baseline and is superseded whenever
    valid foot markers are present in the trial.
\end{itemize}

\section{Reproduction}

\begin{verbatim}
python3 -m pytest tests/unit/motion_matching/test_hip_mirroring.py \
  tests/unit/motion_matching/test_hip_calibration.py \
  tests/unit/motion_matching/pipeline/test_address_feet.py \
  tests/unit/motion_matching/pipeline/test_leg_marker_constraints.py -q
python3 -m scripts.mirror_spec_left_hip \
  docs/development/full_body_models/full_body_spec_anthro_driver.json \
  docs/development/full_body_models/full_body_spec_anthro_iron7.json
python3 -m scripts.mirror_spec_right_knee \
  docs/development/full_body_models/full_body_spec_anthro_driver.json \
  docs/development/full_body_models/full_body_spec_anthro_iron7.json
python3 -m pytest tests/unit/motion_matching/test_knee_axis_convention.py -q
G=docs/development/full_body_models
python3 -m src.shared.python.motion_matching.pipeline.cli \
  --spec $G/full_body_spec_anthro_driver.json --static-seeds \
  --capture driver --out runs/anthro_driver_seeds
python3 -m src.shared.python.motion_matching.pipeline.cli \
  --spec $G/full_body_spec_anthro_iron7.json --static-seeds --zmp-filter \
  --capture iron --out runs/anthro_iron_seeds_zmp
python3 -m src.shared.python.motion_matching.pipeline.cli --address-only \
  --spec $G/full_body_spec_anthro_driver.json --static-seeds \
  --capture driver --engine mujoco --foot-progression capture \
  --out runs/addr_driver_mujoco
python3 -m pytest tests/unit/motion_matching/test_pelvis_alignment_from_spec.py \
  tests/unit/scripts/test_diagnose_address_leg_yaw.py -q
python3 -m scripts.diagnose_address_leg_yaw --capture driver \
  --spec $G/full_body_spec_anthro_driver.json --out runs/legyaw_driver
python3 -m scripts.diagnose_address_leg_yaw --capture iron \
  --spec $G/full_body_spec_anthro_iron7.json --out runs/legyaw_iron
python3 -m scripts.opensim_address_foot_progression \
  --out $G/evidence/foot_progression/opensim_address.json
# Section 8 (#12052); run myosuite with the MyoSuite environment's python.
python3 -m pytest tests/unit/engines/myosuite/test_hip_retarget.py -q
E=$G/evidence/foot_progression/osv6_myosuite/driver   # and .../iron
python3 -m scripts.address_foot_progression_engines evaluate --run-dir $E \
  --label driver --engines mujoco pinocchio --out $E/engines.json
myo-python -m scripts.address_foot_progression_engines evaluate --run-dir $E \
  --label driver --engines myosuite --out $E/engines.json
python3 -m scripts.address_foot_progression_engines render \
  --json $E/engines.json --out-dir stills --engines mujoco pinocchio myosuite
# OSV-5: Golfer realism showcase test and render pipeline (#11731)
python3 -m pytest tests/unit/scripts/test_render_golfer_realism_showcase.py -q
python3 -m scripts.render_golfer_realism_showcase \
  --capture capture-A \
  --driver runs/anthro_driver_seeds/bundle.npz \
  --iron7 runs/anthro_iron_seeds_zmp/bundle.npz \
  --engines mujoco,myosuite,drake,pinocchio,opensim \
  --views face_on,dtl,hands_closeup \
  --speeds 1,0.5 --fps 60 \
  --out ~/Videos/Parity\ Audit/golfer_realism/
\end{verbatim}

\end{document}
